\section{Related Work}
\label{sec:related}

\subsection{Conformal Prediction and Its Sequential Extensions}

Split conformal prediction~\cite{VovkGammermanShafer2005,Lei2018JASA} calibrates a nonconformity score on a held-out set and yields prediction regions with finite-sample, distribution-free marginal coverage regardless of the underlying model. \texttt{calibration.py} follows this standard recipe directly: nonconformity restricted to matched, ground-truth-occupied grid cells, and the usual $\lceil (n+1)(1-\alpha) \rceil / n$ empirical quantile. None of this is novel; it is the substrate the rest of the paper builds on.

That guarantee is single-shot and batch, however, and applying it repeatedly to a stream reintroduces the multiple-testing problem anytime-valid inference exists to solve. Ville's inequality bounds a nonnegative supermartingale's running maximum, and testing by betting~\cite{WaudbySmithRamdas2024JRSSB,Ramdas2023GameTheoretic} recasts a sequential test as a betting fraction against that martingale structure: under the null, wealth cannot systematically grow, and the first crossing of $1/\delta$ is a valid stopping rule at level $\delta$, with no fixed horizon and no correction for how often the test is checked. This is the engine our method runs on (\texttt{betting.py}'s \texttt{WealthProcess}), applied rather than derived.

\subsection{Conformal Prediction for Object Detection}

Two recent lines apply conformal prediction directly to detector outputs, as we do, but in the batch rather than sequential regime. Ries, Sbeyti, Bianco, and Klein~\cite{Ries2026ProbabilisticOD} conformalize bounding-box corners coordinate-wise with a Bonferroni correction, then scale intervals using a probabilistic detector's own aleatoric-uncertainty estimates, reporting sharper regions than unscaled conformal prediction under distribution shift; our nonconformity score is deliberately simpler (an unscaled $\ell_1$ residual) since our contribution is what happens \emph{after} calibration, and their scaling is a natural drop-in improvement to \texttt{calibration.py} if tighter regions matter more than monitoring speed. Kumar, Darabi, Tayebati, and Trivedi's dual-threshold framework~\cite{Kumar2025DualThreshold} pairs a conformal threshold with a ROC-optimized abstention threshold across camera and LiDAR modalities; like~\cite{Geng2025StateDependent,Kumar2025ContextAware}, this is per-prediction batch calibration, answering "is this prediction trustworthy" rather than "has the stream's regime shifted" -- our e-process's question, with an anytime-valid guarantee.

\subsection{Anytime-Valid Failure Detection for Streaming Vision}

The paper closest to ours, and the one we differentiate against explicitly throughout, is Monroy Mu\~noz, Verma, and Timans's sequential object-tracking-failure detector~\cite{MonroyMunoz2026WACV}: an e-process over a bounded per-frame metric $M_t$, with $\lambda_t$ chosen adaptively from $M_t$'s own history via aGRAPA or SF-OGD -- both reproduced as our covariate-blind baselines (\texttt{AGRAPABettor}, \texttt{SFOGDBettor}) so our comparison is against their actual method, not a strawman. Their evaluation spans four 2D visual-tracking benchmarks and is explicitly single-object, per-video; multi-object extension is listed as their own future work. Our spatial multiplicity-control grid is a direct response to that stated gap.

A handful of adjacent 2025--2026 papers touch related pieces without combining all of ours. Geng, Waite, Turnquist, Ivanov, and Ruchkin condition conformal error bounds on a robotic system's dynamical state via reachability analysis~\cite{Geng2025StateDependent} -- covariate-conditioned, but batch, with no e-process and no weather treatment. Kumar, Tayebati, Migliarba, Krishnan, and Trivedi learn context-conditioned nonconformity scores from learned task representations~\cite{Kumar2025ContextAware} -- again batch, and the covariate is a learned embedding rather than a physically interpretable disagreement signal. Liu et al.'s pre-fusion contextual calibration~\cite{Liu2026ValueGated} uses cross-modal agreement to modulate fusion, architecturally close to our monitored detector's own consistency probe, but as a learned gate rather than a statistical covariate, and targets audio-visual tasks. Antonante, Spivak, and Carlone's perception-monitoring work~\cite{Antonante2020Monitoring} frames monitoring graph-theoretically -- the right citation for the general framing, not a competing statistical method. None of these combines (a) anytime-valid sequential testing, (b) multi-object spatial resolution with multiplicity control, and (c) real adverse-weather distribution shift in multimodal 3D driving perception; \cite{MonroyMunoz2026WACV} combines the first without the other two. We also explored conditioning the betting rate on an external covariate -- the same disagreement signal Liu et al.'s gate and our detector's consistency probe both compute -- as a fourth possible axis; Section~\ref{sec:experiments} reports why that attempt did not hold up on real data, so we do not count it among our differentiators.

\subsection{Online Multiplicity Control}

Testing many BEV cells simultaneously, every frame, without inflating the false-discovery rate is a studied problem once the statistics are e-values. Wang and Ramdas's e-BH~\cite{WangRamdas2022JRSSB} rejects hypotheses whose e-values exceed a threshold set by the largest $k$ with $k$-th order statistic $\ge n/(kq)$, controlling FDR under arbitrary dependence -- important here, since neighboring cells' wealth processes are not independent. \texttt{spatial.py} implements this (\texttt{e\_bh\_reject}) alongside a Bonferroni baseline: valid but conservative, versus e-BH's less conservative online alternative. Two 2026 papers extend plain e-BH further. Kuang, Gang, and Xia's SCORE~\cite{Kuang2026SCORE} recovers evidence wasted when an e-value overshoots its threshold, yielding uniformly higher power at the same FDR guarantee -- applied to our grid, a strongly-flagged cell's excess evidence could help a marginal neighbor alarm sooner. Xu, Fischer, and Ramdas~\cite{Xu2026EClosure} generalize closure testing to an online e-closure with compound e-values at $O(\log t)$ per-decision cost, cheap enough to be a plausible replacement for \texttt{e\_bh\_reject} if the real multi-object evaluation we flag as future work shows plain e-BH's power is a bottleneck.

\subsection{Adverse-Weather Perception Data and Multimodal Fusion}

That weather-induced degradation is worth detecting at all, not just engineered away, is corroborated by Vassilev et al.~\cite{Vassilev2026Sensitivity}: benchmarking five detectors, they find diminished lighting and weather-plus-occlusion measurably lengthen a vehicle's effective stopping distance -- our problem exists because that degradation is silent to the stack, not because detectors cannot be made somewhat more robust. Snow-specific datasets exist -- the Canadian Adverse Driving Conditions dataset~\cite{Pitropov2021CADC} and Bijelic et al.'s fog/rain fusion work~\cite{Bijelic2020SeeingThroughFog} -- but neither contains a real in-dataset clear-to-snow transition or supports our calibration/monitoring split. Snowy Scenes~\cite{SnowyScenes2026} (Section~\ref{sec:experiments}) is a more recent (2026) dataset captured entirely in real snow; the price of that specificity is no clear-weather split at all, which we address via a real severity ordering across its own weather categories. The detector we monitor fuses camera and LiDAR BEV features through a Cross-modal Consistency Probe trained under an adversarial consistency objective to detect cross-modal disagreement -- the same signal this paper repurposes as a betting covariate. Section~\ref{sec:method} describes what that probe computes; Section~\ref{sec:experiments} describes what happened when its training objective did not, initially, give it any real signal at all.
